\title{Group Sequence Policy Optimization}

\begin{abstract}
We present Group Sequence Policy Optimization (GSPO), a robust and efficient reinforcement learning method specifically designed for large language model training. Departing from earlier approaches that utilize token-wise importance ratios, GSPO instead calculates the importance ratio from the overall sequence probability and applies clipping, rewards, and optimization at the sequence level. Through empirical analysis, we show that GSPO surpasses the GRPO algorithm in both training efficiency and performance, provides greater stability for Mixture-of-Experts (MoE) RL training scenarios, and offers opportunities to streamline RL system architecture. These strengths have played a significant role in enhancing the latest Qwen3 models.
\end{abstract}

\section{Introduction}
\label{sec:intro}

Reinforcement learning (RL) has become a crucial framework for advancing the development of language models \citep{o1,dpsk-r1,qwq32b,qwen3}. When applied at scale, RL enables language models to address complex challenges, including tasks like competition-level mathematics and programming, by facilitating more extensive and nuanced reasoning over longer sequences.

In order to effectively leverage increased computational resources for scaling RL, it is crucial to ensure that training dynamics remain both stable and resilient. Nevertheless, leading RL methods such as GRPO \citep{grpo} suffer from pronounced instability when used to train very large language models, frequently culminating in significant and unrecoverable collapses of the model \citep{qwen3, minimax-m1}. Such instability poses a significant barrier to further advancing the performance limits of language models via ongoing RL-based training.

In this study, we demonstrate that GRPO's instability arises primarily from an improper use and eventual breakdown of importance sampling weights within its algorithm. This flaw leads to substantial variance in training noise, which grows with longer response sequences and is exacerbated by the clipping mechanism, thereby hastening the collapse of the model.

In response to these fundamental challenges, we introduce \textbf{Group Sequence Policy Optimization (GSPO)}, a novel reinforcement learning method tailored for the training of large language models. The principal contribution of GSPO is its theoretically motivated formulation of the importance ratio via sequence likelihood \citep{click}, adhering closely to the foundation of importance sampling. Furthermore, GSPO evaluates normalized rewards as the relative advantages among several responses to a single query, thereby facilitating a closer correspondence between rewards assigned at the sequence level and the optimization objective.

Our experimental results clearly show that GSPO outperforms GRPO with regard to training stability, efficiency, and overall performance. Notably, GSPO intrinsically addresses the stability issues that typically arise in reinforcement learning for large Mixture-of-Experts (MoE) architectures, obviating the necessity for intricate stabilization methods and paving the way for more streamlined RL systems. The advantages of GSPO have played a pivotal role in achieving notable performance gains in the recent Qwen3 models. We anticipate that GSPO will serve as a reliable and scalable basis for further progress in large-scale reinforcement learning with language models.

\section{Preliminaries}

\paragraph{Notation}
Throughout this work, we represent an autoregressive language model with parameters $\theta$ as a policy $\pi_\theta$. The symbol $x$ stands for a query, while $\mathcal{D}$ refers to the entire set of queries. For a given query $x$ and its response $y$, the conditional probability assigned by the policy $\pi_\theta$ is written as $\pi_\theta (y | x) = \prod_{t=1}^{|y|} \pi_\theta (y_t | x, y_{<t} )$, with $|y|$ indicating the total number of tokens in $y$. A verifier $r$ can be used to assign a reward to the pair $(x, y)$, expressed as $r(x, y) \in [0, 1]$.

\paragraph{Proximal Policy Optimization (PPO)} Given trajectories sampled from the previous policy $\pi_{\theta_\text{old}}$, PPO \citep{ppo} limits the extent of policy modifications to remain close to the previous iteration by introducing a clipping method. In particular, PPO utilizes the following surrogate objective during policy optimization (for clarity, we omit the KL divergence term from further discussion, since it is not central to this work):

\begin{align}
\mathcal{J}_\text{PPO}(\theta) = \mathbb{E}_{ x \sim \mathcal{D},\, y \sim \pi_{\theta_\text{old}}( \cdot | x) }
\left[ \frac{1}{|y|} \sum_{t=1}^{|y|} 
\min \left( w_{t}(\theta) \widehat{A}_{t},  \, \mathrm{clip} \left( w_{t}(\theta), 1 - {\varepsilon}, 1 + {\varepsilon}\right) \widehat{A}_{t} \right)
\right],
\end{align}

Here, the token importance ratio $w_{t}(\theta)$ is given by $
w_{t}(\theta) = \frac{ \pi_{\theta} (y_{t} | x, y_{<t}) }{ \pi_{\theta_\text{old}} (y_{t} | x, y_{<t})} $, with the advantage $\widehat{A}_{t}$ for $y_t$ being approximated through a separate value function, and $\varepsilon$ indicating the interval for clipping these importance ratios.

A central difficulty encountered when implementing PPO is its substantial dependence on the value model. In most cases, the value model matches the policy model in scale, resulting in significant demands on both memory and computational resources. Additionally, the overall performance of the algorithm depends critically on the accuracy of the value estimation. Not only is obtaining a trustworthy value model a complex task, but extending its reliability to accommodate longer outputs and increasingly intricate problems further exacerbates this issue.

\paragraph{Group Relative Policy Optimization (GRPO)} GRPO \citep{grpo} eliminates the dependence on a value model by determining the comparative advantage for each answer among a collection of responses generated for the same prompt. In particular, GRPO seeks to maximize the following objective:

\begin{align}
\label{equ:grpo}
\mathcal{J}_\text{GRPO}(\theta) = \mathbb{E}_{ x \sim \mathcal{D},\, \{y_i\}_{i=1}^G \sim \pi_{\theta_\text{old}}( \cdot | x) }
\left[ \frac{1}{G} \sum_{i=1}^{G} \frac{1}{|y_i|} \sum_{t=1}^{|y_i|} 
\min \left( w_{i,t}(\theta) \widehat{A}_{i,t},  \, \mathrm{clip} \left( w_{i,t}(\theta), 1 - {\varepsilon}, 1 + {\varepsilon}\right) \widehat{A}_{i,t} \right)
\right],
\end{align}

Here, $G$ denotes the number of responses generated for each input $x$ (referred to as the group size), and for token $y_{i,t}$, the terms $w_{i,t}(\theta)$ and $\widehat{A}_{i,t}$ represent the importance ratio and the estimated advantage, respectively:

\begin{align}
    w_{i,t}(\theta)=\frac{ \pi_{\theta} (y_{i,t} | x, y_{i,<t}) }{ \pi_{\theta_\text{old}} (y_{i,t} | x,y_{i,<t})},\quad\ 
    \widehat{A}_{i,t} = \widehat{A}_{i} = \frac{r(x, y_i) - \mathrm{mean} \left( \{ r(x, y_i) \}_{i=1}^G \right) }{ \mathrm{std} \left( \{ r(x, y_i) \}_{i=1}^G \right) },
\end{align}

respectively, with every token in $y_i$ being associated with the identical advantage value, $\widehat{A}_{i}$.

\section{Motivation}
\label{sec:motivation}

As models become larger, incorporate greater sparsity (such as seen in Mixture-of-Experts architectures), and produce longer responses, it becomes increasingly important to use substantial rollout batch sizes to ensure efficient utilization of hardware during reinforcement learning. To enhance sample efficiency, it is commonly accepted to subdivide these large batches of rollout data into several mini-batches during gradient update steps. However, this strategy leads naturally to an off-policy learning context, since the sampled responses $y$ are generated from a previous policy $\pi_{\theta_\text{old}}$ rather than the actively updated policy $\pi_\theta$. This situation underpins the rationale for employing the clipping technique in algorithms like PPO and GRPO, as such mechanisms are designed to limit the influence of samples that diverge excessively from the current policy, thereby ensuring stable and reliable gradient estimation.

Although techniques such as clipping seek to address the off-policy mismatch, our analysis reveals a deeper problem inherent in GRPO: \textit{the objective itself lacks a well-defined formulation}. This deficiency is especially pronounced when applying GRPO to train large-scale models for tasks requiring extended outputs, often resulting in severe degradation of the model's performance. The fundamental issue arises from the inappropriate application of importance sampling weights. Importance sampling is traditionally employed to compute the expected value of a function $f$ with respect to a target distribution $\pi_\text{tar}$ by adjusting the weights of samples originally obtained from a different, behavior distribution $\pi_\text{beh}$:

\begin{align}
\mathbb{E}_{z \sim \pi_\text{tar}} \left[ f(z) \right] = \mathbb{E}_{z \sim \pi_\text{beh}} \left[ \frac{ \pi_\text{tar} (z) }{ \pi_\text{beh} (z)} \, f(z) \right].
\end{align}

Importantly, this approach depends on taking an average over a large number of samples ($N \gg 1$) drawn from the behavior distribution $\pi_\text{beh}$, which allows the importance weight $\frac{ \pi_\text{tar} (z) }{ \pi_\text{beh} (z)}$ to adequately adjust for differences between the distributions.

Conversely, GRPO incorporates the importance weight $\frac{ \pi_{\theta} (y_{i,t} | x, y_{i,<t}) }{ \pi_{\theta_\text{old}} (y_{i,t} | x,y_{i,<t})}$ at every token step $t$. Because this adjustment relies on just a single draw $y_{i,t}$ from the next-token probability distribution $\pi_{\theta_\text{old}}( \cdot | x, y_{i, <t})$, it does not adequately correct for discrepancies between distributions. Rather, this method injects significant variance into the gradient estimates, which accumulates over longer sequences—a phenomenon further aggravated by the use of clipping. Our empirical results indicate that this process can precipitate model collapse, from which recovery is rarely possible. Once such a collapse has taken place, attempts to continue training—even by reverting to previous checkpoints, finely adjusting hyperparameters like clipping bounds, lengthening the generation, or altering RL sampling procedures—do not restore effective learning.

The preceding analysis highlights an inherent flaw within GRPO’s architecture. The ineffectiveness of token-level importance weighting underscores a crucial concept: \textbf{the granularity of the optimization target must align with that of the reward}. Given that rewards are assigned for complete sequences, implementing off-policy adjustments on a per-token basis is likely inappropriate. Consequently, we are prompted to abandon the token-centric approach in favor of formulating importance weights and conducting optimization at the \textit{sequence level} instead.

\section{Algorithm}

\subsection{GSPO: Group Sequence Policy Optimization}

Although using the token-level importance weight $\frac{ \pi_{\theta} (y_{i,t} | x, y_{i,<t}) }{ \pi_{\theta_\text{old}} (y_{i,t} | x,y_{i,<t})}$ presents challenges within GRPO, we note that the \textit{sequence-level} importance weight $\frac{ \pi_{\theta} (y | x) }{ \pi_{\theta_\text{old}} (y | x)}$ in the context of language generation holds clear theoretical significance. Specifically, it quantifies the extent to which a sampled response $y$ from $\pi_{\theta_\text{old}} (\cdot | x)$ diverges from the updated distribution $\pi_{\theta} (\cdot | x)$. This correspondence is consistent with the reward defined at the sequence level and provides an interpretable metric for the operation of the clipping mechanism.

Drawing from this basic insight, we introduce the \textbf{Group Sequence Policy Optimization (GSPO)} algorithm. The GSPO framework utilizes the subsequent objective for optimization at the sequence level:

\begin{align}
\mathcal{J}_\text{GSPO} (\theta) =
\mathbb{E}_{ x \sim \mathcal{D},\, \{y_i\}_{i=1}^G \sim \pi_{\theta_\text{old}}( \cdot | x) }
\left[ 
\frac{1}{G} \sum_{i=1}^{G}
\min \left( s_{i}(\theta)  \widehat{A}_{i},  \, \mathrm{clip} \left( s_{i}(\theta), 1 - {\varepsilon}, 1 + {\varepsilon} \right) \widehat{A}_{i} \right) 
\right],
\label{equ:gspo}
\end{align}

where we utilize the estimation of advantages based on groupings:

\begin{align}
\widehat{A}_{i} = \frac{r(x, y_i) - \mathrm{mean} \left( \{ r(x, y_i) \}_{i=1}^G \right) }{ \mathrm{std} \left( \{ r(x, y_i) \}_{i=1}^G \right) },
\end{align}

and utilize sequence likelihood to establish the importance ratio $s_{i}(\theta)$, as in \citet{click}:

\begin{align}
s_{i}(\theta) = \left( \frac{ \pi_{\theta} (y_i | x) }{ \pi_{\theta_\text{old}} (y_i | x)} \right)^{\frac{1}{|y_i|}}
=
\exp \left( \frac{1}{|y_i|} \sum_{t=1}^{|y_i|} \log \frac{ \pi_{\theta} (y_{i,t} | x, y_{i,<t}) }{ \pi_{\theta_\text{old}} (y_{i,t} | x,y_{i,<t})} \right).
\end{align}

Consequently, GSPO implements clipping at the level of full responses rather than at individual tokens, effectively filtering out highly ``off-policy'' samples during gradient calculation. This approach aligns with the objectives of sequence-level reward assignment and the overall optimization process. Additionally, we utilize length normalization in $s_{i}(\theta)$, aiming to decrease its variance and standardize $s_{i}(\theta)$ within a consistent numerical interval. Without such normalization, changes in the likelihoods of just a few tokens could induce significant variability in the sequence-level importance ratio, necessitating distinct clipping thresholds for responses of varying lengths. It is also important to highlight that, due to differing formulations of the importance ratios, the clipping thresholds used by GSPO often differ in scale compared to those in preceding methods such as GRPO.

\subsection{Gradient Analysis}
\label{subsec:gradient}

The gradient of the GSPO objective can be obtained through the following derivation (for simplicity, clipping is not included):

\begin{align}
\nabla_{\theta} \mathcal{J}_\text{GSPO} (\theta)
=&\ 
\nabla_{\theta} \mathbb{E}_{ x \sim \mathcal{D},\, \{y_i\}_{i=1}^G \sim \pi_{\theta_\text{old}}( \cdot | x) }
\left[ 
\frac{1}{G} \sum_{i=1}^{G} s_{i}(\theta) \widehat{A}_{i}
\right] \\
= &\ 
\mathbb{E}_{ x \sim \mathcal{D},\, \{y_i\}_{i=1}^G \sim \pi_{\theta_\text{old}}( \cdot | x) }
\left[ 
\frac{1}{G} \sum_{i=1}^{G}
s_{i}(\theta) \widehat{A}_{i}
\cdot \nabla_{\theta} \log s_{i}(\theta)
\right] \\
=&\ 
\mathbb{E}_{ x \sim \mathcal{D},\, \{y_i\}_{i=1}^G \sim \pi_{\theta_\text{old}}( \cdot | x) }
\left[ 
\frac{1}{G} \sum_{i=1}^{G} 
\left( \frac{ \pi_{\theta} (y_i | x) }{ \pi_{\theta_\text{old}} (y_i | x)} \right)^{\frac{1}{|y_i|}} \widehat{A}_{i} 
\cdot \frac{1}{|y_i|} \sum_{t=1}^{|y_i|} \nabla_{\theta} \log \pi_{\theta} (y_{i,t} | x, y_{i,<t}) 
\right].
\label{equ:gspo_grad}
\end{align}

For reference, the gradient associated with the GRPO objective can be expressed as follows (where $\widehat{A}_{i,t}$ simplifies to $\widehat{A}_{i}$):

\begin{align}
\nabla_{\theta} \mathcal{J}_\text{GRPO} (\theta) 
=&\ 
\nabla_{\theta} \mathbb{E}_{ x \sim \mathcal{D},\, \{y_i\}_{i=1}^G \sim \pi_{\theta_\text{old}}( \cdot | x) }
\left[ \frac{1}{G} \sum_{i=1}^{G} \frac{1}{|y_i|} \sum_{t=1}^{|y_i|} 
w_{i,t}(\theta) \widehat{A}_{i,t}
\right] \\
=&\
\mathbb{E}_{ x \sim \mathcal{D},\, \{y_i\}_{i=1}^G \sim \pi_{\theta_\text{old}}( \cdot | x) }
\left[ \frac{1}{G} \sum_{i=1}^{G} \widehat{A}_{i} 
\cdot \frac{1}{|y_i|} \sum_{t=1}^{|y_i|} 
\frac{ \pi_{\theta} (y_{i,t} | x, y_{i,<t}) }{ \pi_{\theta_\text{old}} (y_{i,t} | x,y_{i,<t})} 
\nabla_{\theta} \log \pi_{\theta} (y_{i,t} | x, y_{i,<t})  
\right].
\end{align}

As a result, the primary difference between GSPO and GRPO concerns \textit{the manner in which each method assigns weights to the gradients of the token log likelihoods}. GRPO applies a weighting scheme based on the ``importance weight'' $\frac{ \pi_{\theta} (y_{i,t} | x, y_{i,<t}) }{ \pi_{\theta_\text{old}} (y_{i,t} | x,y_{i,<t})}$ for each token, causing the influence of individual tokens to depend on these potentially disparate weights. These weights may range within $(0, 1+\varepsilon]$ (when $\widehat{A}_{i} > 0$) or $[1-\varepsilon, +\infty)$ (when $\widehat{A}_{i} < 0$), and their variability can have significant and accumulating effects throughout the course of training, possibly resulting in unstable learning behavior. By contrast, GSPO ensures uniform weighting across all tokens within a response, thereby removing the source of instability associated with GRPO's non-uniform weighting.

\subsection{GSPO-token: A Token-level Objective Variant}

In situations such as multi-turn RL, it can be beneficial to perform advantage adjustment at a more granular level than the full sequence. Therefore, we propose a token-level extension of GSPO, referred to as \textbf{GSPO-token}, which enables the assignment of token-specific advantages:

\begin{align}
\mathcal{J}_\text{GSPO-token}(\theta)
=
\mathbb{E}_{ x \sim \mathcal{D},\, \{y_i\}_{i=1}^G \sim \pi_{\theta_\text{old}}( \cdot | x) }
\left[ 
\frac{1}{G} \sum_{i=1}^{G} \frac{1}{|y_i|} \sum_{t=1}^{|y_i|} 
\min \left( s_{i,t}(\theta) \widehat{A}_{i,t},  \, \mathrm{clip} \left( s_{i,t}(\theta), 1 - {\varepsilon}, 1 + {\varepsilon} \right) \widehat{A}_{i,t} \right) 
\right],
\label{equ:gspo-token}
\end{align}

in which

\begin{align}
s_{i,t}(\theta) 
= \mathrm{sg} \left[ s_{i}(\theta) \right]  \cdot \frac{ \pi_{\theta} (y_{i,t} | x, y_{i,<t}) }{ \mathrm{sg} \left[ \pi_{\theta} (y_{i,t} | x, y_{i,<t}) \right] },
\end{align}

and $\mathrm{sg}[\cdot]$ indicates extracting the value without allowing gradients to propagate, aligning with the \texttt{detach} function in PyTorch. The derivative of GSPO-token with respect to its parameters can be formulated as follows:

\begin{align}
\nabla_{\theta} \mathcal{J}_\text{GSPO-token}(\theta)
=&\ 
\nabla_{\theta} \mathbb{E}_{ x \sim \mathcal{D},\, \{y_i\}_{i=1}^G \sim \pi_{\theta_\text{old}}( \cdot | x) }
\left[ 
\frac{1}{G} \sum_{i=1}^{G}
\frac{1}{|y_i|} \sum_{t=1}^{|y_i|} 
s_{i,t}(\theta) \widehat{A}_{i,t}
\right] \\
=&\ 
\mathbb{E}_{ x \sim \mathcal{D},\, \{y_i\}_{i=1}^G \sim \pi_{\theta_\text{old}}( \cdot | x) }
\left[ 
\frac{1}{G} \sum_{i=1}^{G} s_i (\theta) \cdot \frac{1}{|y_i|} \sum_{t=1}^{|y_i|} 
\widehat{A}_{i,t} \frac{ \nabla_{\theta} \pi_{\theta} (y_{i,t} | x, y_{i,<t}) }{ \pi_{\theta} (y_{i,t} | x, y_{i,<t}) }
\right] \\
=&\ 
\mathbb{E}_{ x \sim \mathcal{D},\, \{y_i\}_{i=1}^G \sim \pi_{\theta_\text{old}}( \cdot | x) }
\left[ 
\frac{1}{G} \sum_{i=1}^{G}  \left( \frac{ \pi_{\theta} (y_i | x) }{ \pi_{\theta_\text{old}} (y_i | x)} \right)^{\frac{1}{|y_i|}}
\cdot \frac{1}{|y_i|} \sum_{t=1}^{|y_i|}
\widehat{A}_{i,t} \nabla_{\theta} \log \pi_{\theta} (y_{i,t} | x, y_{i,<t})  
\right].
\label{equ:gspo-token_grad}
\end{align}

Observe that the expression $\frac{ \pi_{\theta} (y_{i,t} | x, y_{i,<t}) }{ \mathrm{sg} \left[ \pi_{\theta} (y_{i,t} | x, y_{i,<t}) \right] }$ always evaluates to 1, which means that $s_{i,t}(\theta)$ is, in practice, equal to $s_{i}(\theta)$. By analyzing Equation~\eqref{equ:gspo} in relation to~\eqref{equ:gspo-token}, and Equation~\eqref{equ:gspo_grad} relative to~\eqref{equ:gspo-token_grad}, we can see that GSPO and GSPO-token are functionally equivalent—both in terms of the objective, the criteria for clipping, and their gradients—so long as the advantage values assigned to all tokens in the response $y_i$ are held constant (i.e., $\widehat{A}_{i,t} = \widehat{A}_{i}$). Nevertheless, GSPO-token offers greater adaptability, as it allows for the assignment of distinct advantages at the level of individual tokens.

\section{Experiments and Discussion}

\subsection{Empirical Results}

Our experiments utilize a cold-start model, initialized from Qwen3-30B-A3B-Base, for fine-tuning. We present both the progression of training rewards and the evolution of model performance metrics on several benchmarks: AIME'24 (mean Pass@1 over 32 samplings), LiveCodeBench (202410-202502, mean Pass@1 across 8 samplings), and CodeForces (measured by Elo Rating). For reinforcement learning training, each rollout batch is divided into four mini-batches for the purpose of gradient computation and updates. When applying GSPO, the clipping bounds in Equation~\eqref{equ:gspo} are chosen as 3e-4 (left) and 4e-4 (right). Our comparison employs GRPO as a baseline, with its clipping intervals in Equation~\eqref{equ:grpo} set to 0.2 and 0.27, respectively; these values are selected following extensive tuning to guarantee an equitable evaluation. Importantly, we emphasize that successful training convergence for MoE RL with GRPO depends on employing the Routing Replay strategy, which we further elaborate upon in \S~\ref{subsec:routing_replay}. In contrast, \textit{GSPO removes the necessity for this technique}.

As depicted in Figure~\ref{fig:results}, the training process employing GSPO remains stable at all stages. Notably, \textbf{GSPO facilitates ongoing performance gains by scaling up training compute, periodically refreshing the query set, and increasing the generation length}. In addition, GSPO exhibits greater training efficiency than GRPO, obtaining superior accuracy and benchmark results while utilizing the same training compute and number of queries. Lastly, GSPO has been effectively deployed for RL training on the most recent Qwen3 models, providing compelling evidence of its ability to harness RL scaling for large language models.

\subsection{Curious Observation on Clipping Fractions}

One fundamental difference between GSPO and GRPO lies in the scope of their clipping strategies: GSPO operates at the sequence (response) level, whereas GRPO applies clipping at the individual token level. As illustrated in Figure~\ref{fig:clipping}, this leads to a disparity of approximately two orders of magnitude in the proportions of clipped tokens between the two methods, a gap that persists regardless of changes to the clipping thresholds. Nevertheless, even though GSPO discards a much larger proportion of tokens—which means that fewer are available for model training or for estimating gradients—it consistently achieves better training efficiency relative to GRPO. This seemingly paradoxical result, where more extensive token clipping correlates with improved efficiency, suggests that GRPO’s approach of deriving gradients from single tokens introduces substantial noise and hampers effective exploitation of the available samples. By focusing on the response level, GSPO instead delivers a more stable and informative training signal.

\subsection{Benefit of GSPO for MoE Training}
\label{subsec:routing_replay}

\paragraph{Background}
Unlike the RL training of dense models, MoE models, due to their sparse activation patterns, present distinctive challenges related to stability. Notably, our investigation revealed that when utilizing the GRPO algorithm, the inherent \textit{expert-activation volatility} in MoE models can obstruct the effective convergence of RL training. Specifically, following one or multiple gradient update steps, the set of experts engaged for an identical response may shift considerably. For instance, in experiments with the 48-layer Qwen3-30B-A3B-Base architecture, analysis showed that, for the same rollout sample after each RL gradient update, approximately 10\% of the experts selected under the updated policy $\pi_{\theta}$ differ from those chosen under the previous policy $\pi_{\theta_\text{old}}$. This effect is especially pronounced in MoE models with greater depth, leading to large variations in the token-level importance ratios $w_{i,t}(\theta) = \frac{ \pi_{\theta} (y_{i,t} | x, y_{i,<t}) }{ \pi_{\theta_\text{old}} (y_{i,t} | x,y_{i,<t})}$ and ultimately undermining their reliability, as explored in \S~\ref{sec:motivation} and~\ref{subsec:gradient}. As a result, this instability disrupts the expected progression of RL training convergence.

\paragraph{Our Previous Approach} In addressing this issue, our prior work utilized the \textbf{Routing Replay} training methodology. This method involves storing the expert selections made by $\pi_{\theta_\text{old}}$ and subsequently reenacting these same routing decisions within $\pi_{\theta}$ during the calculation of the importance weights $w_{i,t}(\theta) = \frac{ \pi_{\theta} (y_{i,t} | x, y_{i,<t}) }{ \pi_{\theta_\text{old}} (y_{i,t} | x,y_{i,<t})}$. Consequently, for each individual token $y_{i,t}$, both $\pi_{\theta} (y_{i,t} | x, y_{i,<t})$ and $\pi_{\theta_\text{old}} (y_{i,t} | x,y_{i,<t})$ make use of an identical activated subnetwork, thereby promoting stable token-level importance ratios. This also guarantees that the gradient-based updates optimize the same subset of activated experts. As illustrated in Figure~\ref{fig:routing_replay}, Routing Replay emerges as a vital component, facilitating reliable convergence in GRPO training when applied to MoE models.

\paragraph{Benefit of GSPO} While Routing Replay allows for successful GRPO training of MoE models, its reliance on the reuse of routing decisions introduces extra burdens in terms of memory usage and communication, and it can restrict the effective utilization of model capacity. By contrast, as illustrated in Figure~\ref{fig:results}, GSPO removes the necessity for Routing Replay, making it possible to compute the importance ratios $s_i(\theta)$ directly and in a standard manner, ensuring stable and regular convergence. The central observation is that GSPO prioritizes the sequence-level likelihood (specifically, $\pi_{\theta} (y_i | x)$), rather than the likelihood of individual tokens (namely, $\pi_{\theta} (y_{i,t} | x, y_{i,<t})$), making it robust against minor token-level variations. Given that the MoE model consistently retains its overall language modeling ability, the likelihood at the sequence level remains comparatively stable. In essence, GSPO effectively addresses the instability in expert activation common to MoE models, thereby eliminating the necessity for more elaborate solutions such as Routing Replay. This results in a more straightforward and dependable training procedure, and enables the model to exploit its full modeling capacity without imposed restrictions.

\subsection{Benefit of GSPO for RL Infrastructure}

Due to the variations in numerical precision between training engines (such as Megatron) and inference engines (like SGLang and vLLM), it is common practice to recalculate the likelihoods of generated responses from the old policy $\pi_{\theta_\text{old}}$ using the training engine. In contrast, GSPO leverages sequence-level likelihoods during optimization rather than token-level ones, and sequence-level likelihoods are generally more robust to minor precision mismatches. As a result, GSPO enables practitioners to utilize the likelihoods provided directly by the inference engine for optimization purposes, eliminating the requirement to recompute them with the training engine. This feature is particularly advantageous in cases such as partial rollout, multi-turn RL, and architectures that separate training and inference workloads.

\section{Conclusion}

We introduce Group Sequence Policy Optimization (GSPO), an innovative reinforcement learning algorithm tailored for training large language models. GSPO employs importance sampling as its core principle, formulating importance ratios with respect to sequence likelihoods and applying sequence-level strategies such as clipping, rewarding, and optimization. In empirical evaluations, GSPO delivers substantially enhanced training stability, efficiency, and effectiveness relative to GRPO, and proves especially advantageous for large-scale reinforcement learning in MoE architectures. These characteristics underpin the significant advancements observed in the latest Qwen3 models. As GSPO establishes itself as a scalable foundation, we anticipate further expanding reinforcement learning and achieving meaningful breakthroughs in artificial intelligence.

\end{document}